\subsubsection{A rational tail that cannot be filled in}

\begin{example}[Refuted: rational exceptional tails may remain wildly ramified]
\label{D:pure-root-rational-tail-obstruction}
Rationality of the tails introduced by a purely inseparable
modification does not imply that a pulled-back cover extends
\'etale across them.  The following local example proves
this failure even when the normalized pure cover is smooth.
It is not a counterexample to D.1 or D.2(ii).
\end{example}

\begin{proof}
Let $k$ be algebraically closed of characteristic two and put
\[
       Y=\operatorname{Spec}k[x,t],\qquad C=(t=0).
\]
Adjoin $z$ with $z^2=t+x^3$.  This is a purely inseparable
field extension of degree two: $t+x^3$ is not a square in
$k(x,t)$, for example because its $t$-derivative is one.
Its normalization over $Y$ is
\[
       Y'=\operatorname{Spec}k[x,z],\qquad t=z^2+x^3.
\]
Indeed this polynomial ring is finite over $k[x,t]$,
normal, and has the required fraction field.  Thus $Y'$
is smooth and locally factorial.

The inverse image $C'$ of $C$ is the cusp $z^2=x^3$.
Its normalization has coordinates
\[
                      x=w^2,\qquad z=w^3.
\]
At the cusp the element $w=z/x$ is not in
$k[[w^2,w^3]]$, whose nonconstant terms have order at
least two.  Consequently there is no morphism
$C'\to(C')^{\rm nor}$ inverse to normalization.
In particular no retraction of the formal neighborhood
onto this smooth normalized main curve can restrict
to such an inverse on the reduced curve.  Retaining
the locally factorial surface $Y'$ does not resolve
this problem with its main curve.

Blow up the origin of $Y'$ and use the chart
$x=s$, $z=sy$.  On this chart
\[
                  t=s^2(y^2+s).
 \tag{PT}
\]
The exceptional divisor $R=(s=0)$ is rational.  The
strict transform $E$ of $C'$ is $y^2+s=0$ and meets
$R$ only at $y=0$.

On $Y\setminus C$ consider the cover
\[
                         u^2+u=\frac1t.
 \tag{PA}
\]
It is finite \'etale of degree two, since the derivative
with respect to $u$ is one.  It is connected: if
$b^2+b=1/t$ in $k(x,t)$, the $t$-valuation of $b$
would have to be negative, and the left side would
then have an even negative valuation, whereas the
right side has valuation $-1$.  This is also the
usual Artin--Schreier torsor construction
\cite[section 59.63, tag 0A3J]{StacksBFiniteCone}.

Pull this cover back along the pure normalization and
the blowup.  Use (PT) and put
$v=u+1/(sy)$.  In characteristic two the exact
rational-function identity is
\[
 v^2+v
 =\frac{1}{s^2(y^2+s)}
       +\frac1{s^2y^2}+\frac1{sy}
 =\frac{1+y^3+sy}{s\,y^2(y^2+s)}.                    \tag{PW}
\]
At the generic point of $R$, the right side has a
simple pole, with nonzero leading coefficient
$(1+y^3)/y^4$ in $k(y)$.  In any extension of this
DVR valuation, normalized to give $s$ value one,
equation (PW) forces $v$ to have value $-1/2$.
The ramification index is therefore two, since the
field extension has degree two.  In particular it
is wild at this rational exceptional divisor.

The generic point of $R$ is outside $E$.  Thus a
cover that is \'etale on the complement of $E\cup R$
in this chart need not extend \'etale to the complement
of $E$ obtained by filling in $R\setminus E$.
Further embedded resolution does not change the
valuation at the generic point of the strict
transform of $R$.

The curve $C$ in this example is affine and is not
the proper exceptional curve of a contracted
surface germ.  The conclusion concerns precisely
the failed extension inference, not the existence
of a global cover violating any finite-generation
statement.  The exact companion program
\texttt{check\_pure\_root\_tail.py} checks thirteen
finite polynomial and rational-function identities;
the ramification and scope conclusions are the
arguments above.
\end{proof}
